% GENERATED FILE. Edit canonical lessons, then run npm run generate:books.
% canonical-source-hash: fnv1a64:cf7a31abb6f6f373
% canonical-lessons: JA-C396-yukisaki, JA-C396-samayou, JA-C396-tadoritsuku, JA-C396-burabura, JA-C396-urouro

\chapter{Where Are You Headed?}
\label{ch:ja-where-are-you-headed}
\hlchaptermodality{396}

\begin{chapteropening}
\textbf{By the end of this chapter:} \emph{I can talk about destinations and wandering about.}

Complete the last lesson of chapter 396: I can talk about destinations and wandering about.
\end{chapteropening}

\section[\texorpdfstring{\ja{ゆきさき} (yukisaki)}{ゆきさき (yukisaki)}]{\ja{ゆきさき} (yukisaki) — a destination}

\label{lesson:JA-C396-yukisaki}



\begin{quote}
Before the new one: say the Japanese for to loosen, then the Japanese for to go past one's turning.
\end{quote}

\subsection*{You'll want to know: \ja{ゆきさき}}
\textbf{\ja{ゆきさき}} — \emph{yukisaki} — \textquotedblleft{}a destination\textquotedblright{}.

\begin{quote}
\textbf{\ja{でんしゃの} \ja{ゆきさきを} \ja{たしかめました。}} — \emph{densha no yukisaki o tashikamemashita} — \textquotedblleft{}I made sure of the train's destination.\textquotedblright{}
\end{quote}

\begin{cousinweb}
\textbf{\ja{ゆきさき}} (\emph{yukisaki}) is a native Japanese word (wago), inherited rather than borrowed from Chinese, built from \emph{yuki}, \textquotedblleft{}going\textquotedblright{} (from \emph{yuku, another form of iku}) and \emph{saki}, \textquotedblleft{}ahead\textquotedblright{}; its kanji were attached later as kun'yomi (JMdict).
\end{cousinweb}

\begin{grammarlens}[title={What you've built}]
One more word for getting around.
\end{grammarlens}

\subsection*{Your turn}
\begin{itemize}
\raggedright
  \item \emph{Say it:} \emph{yukisaki}
  \item \emph{Say it:} \emph{yukisaki}, once more
  \item \emph{Say it:} \emph{densha no yukisaki o tashikamemashita} — I made sure of the train's destination.
  \item \emph{Recall:} say the Japanese for to loosen, then the Japanese for to go past one's turning, then say \emph{yukisaki} again
\end{itemize}

\begin{tcolorbox}[breakable,colback=teal!4,colframe=teal!35!black,title={Before you move on}]
What does \ja{ゆきさき} mean? (\textquotedblleft{}A destination\textquotedblright{}.) Say it once more.
\end{tcolorbox}
\section[\texorpdfstring{\ja{さまよう} (samayou)}{さまよう (samayou)}]{\ja{さまよう} (samayou) — to wander about}

\label{lesson:JA-C396-samayou}



\begin{quote}
Before the new one: say the Japanese for to go past one's turning, then the Japanese for a destination.
\end{quote}

\subsection*{You'll want to know: \ja{さまよう}}
\textbf{\ja{さまよう}} — \emph{samayou} — \textquotedblleft{}to wander about\textquotedblright{}.

\begin{quote}
\textbf{\ja{ちずを} \ja{なくして、まちを} \ja{さまよいました。}} — \emph{chizu o nakushite, machi o samayoimashita} — \textquotedblleft{}I lost my map and wandered about the town.\textquotedblright{}
\end{quote}

Polite \textbf{\ja{さまよいます}}, \emph{samayoimasu}; a godan verb. Past (polite): \textbf{\ja{あにが} \ja{まちを} \ja{さまよいました。}}, \emph{ani ga machi o samayoimashita}.

\begin{cousinweb}
\textbf{\ja{さまよう}} (\emph{samayou}) is a native Japanese word (wago), inherited rather than borrowed from Chinese; its kanji was attached later as a kun'yomi (JMdict).
\end{cousinweb}

\begin{grammarlens}[title={What you've built}]
One more word for getting around.
\end{grammarlens}

\subsection*{Your turn}
\begin{itemize}
\raggedright
  \item \emph{Say it:} \emph{samayou}
  \item \emph{Say it:} \emph{samayou}, once more
  \item \emph{Say it:} \emph{chizu o nakushite, machi o samayoimashita} — I lost my map and wandered about the town.
  \item \emph{Recall:} say the Japanese for to go past one's turning, then the Japanese for a destination, then say \emph{samayou} again
\end{itemize}

\begin{tcolorbox}[breakable,colback=teal!4,colframe=teal!35!black,title={Before you move on}]
What does \ja{さまよう} mean? (\textquotedblleft{}To wander about\textquotedblright{}.) Say it once more.
\end{tcolorbox}
\section[\texorpdfstring{\ja{たどりつく} (tadoritsuku)}{たどりつく (tadoritsuku)}]{\ja{たどりつく} (tadoritsuku) — to finally reach}

\label{lesson:JA-C396-tadoritsuku}



\begin{quote}
Before the new one: say the Japanese for a destination, then the Japanese for to wander about.
\end{quote}

\subsection*{You'll want to know: \ja{たどりつく}}
\textbf{\ja{たどりつく}} — \emph{tadoritsuku} — \textquotedblleft{}to finally reach\textquotedblright{}.

\begin{quote}
\textbf{\ja{ばんに} \ja{やどに} \ja{たどりつきました。}} — \emph{ban ni yado ni tadoritsukimashita} — \textquotedblleft{}We finally reached the inn in the evening.\textquotedblright{}
\end{quote}

Polite \textbf{\ja{たどりつきます}}, \emph{tadoritsukimasu}; a godan verb. Past (polite): \textbf{\ja{ふねが} \ja{しまに} \ja{たどりつきました。}}, \emph{fune ga shima ni tadoritsukimashita}.

\begin{cousinweb}
\textbf{\ja{たどりつく}} (\emph{tadoritsuku}) is a native Japanese word (wago), inherited rather than borrowed from Chinese, built from \emph{tadori}, \textquotedblleft{}following a path\textquotedblright{} (from \emph{tadoru}) and \emph{tsuku}, \textquotedblleft{}to arrive\textquotedblright{}; its kanji were attached later as kun'yomi (JMdict).
\end{cousinweb}

\begin{grammarlens}[title={What you've built}]
One more word for getting around.
\end{grammarlens}

\subsection*{Your turn}
\begin{itemize}
\raggedright
  \item \emph{Say it:} \emph{tadoritsuku}
  \item \emph{Say it:} \emph{tadoritsuku}, once more
  \item \emph{Say it:} \emph{ban ni yado ni tadoritsukimashita} — We finally reached the inn in the evening.
  \item \emph{Recall:} say the Japanese for a destination, then the Japanese for to wander about, then say \emph{tadoritsuku} again
\end{itemize}

\begin{tcolorbox}[breakable,colback=teal!4,colframe=teal!35!black,title={Before you move on}]
What does \ja{たどりつく} mean? (\textquotedblleft{}To finally reach\textquotedblright{}.) Say it once more.
\end{tcolorbox}
\section[\texorpdfstring{\ja{ぶらぶら} (burabura)}{ぶらぶら (burabura)}]{\ja{ぶらぶら} (burabura) — strolling; aimlessly}

\label{lesson:JA-C396-burabura}



\begin{quote}
Before the new one: say the Japanese for to wander about, then the Japanese for to finally reach.
\end{quote}

\subsection*{You'll want to know: \ja{ぶらぶら}}
\textbf{\ja{ぶらぶら}} — \emph{burabura} — \textquotedblleft{}strolling; aimlessly\textquotedblright{}.

\begin{quote}
\textbf{\ja{しゅうまつは} \ja{まちを} \ja{ぶらぶら} \ja{あるきました。}} — \emph{shūmatsu wa machi o burabura arukimashita} — \textquotedblleft{}At the weekend I strolled around the town.\textquotedblright{}
\end{quote}

\begin{cousinweb}
\textbf{\ja{ぶらぶら}} (\emph{burabura}) is a native mimetic word that pictures an easy, aimless stroll: an imitative Japanese form, inherited rather than borrowed, whose older origin is not settled (JMdict).
\end{cousinweb}

\begin{grammarlens}[title={What you've built}]
One more word for getting around.
\end{grammarlens}

\subsection*{Your turn}
\begin{itemize}
\raggedright
  \item \emph{Say it:} \emph{burabura}
  \item \emph{Say it:} \emph{burabura}, once more
  \item \emph{Say it:} \emph{shūmatsu wa machi o burabura arukimashita} — At the weekend I strolled around the town.
  \item \emph{Recall:} say the Japanese for to wander about, then the Japanese for to finally reach, then say \emph{burabura} again
\end{itemize}

\begin{tcolorbox}[breakable,colback=teal!4,colframe=teal!35!black,title={Before you move on}]
What does \ja{ぶらぶら} mean? (\textquotedblleft{}Strolling; aimlessly\textquotedblright{}.) Say it once more.
\end{tcolorbox}
\section[\texorpdfstring{\ja{うろうろ} (urōro)}{うろうろ (urōro)}]{\ja{うろうろ} (urōro) — wandering restlessly}

\label{lesson:JA-C396-urouro}



\begin{quote}
Before the new one: say the Japanese for to finally reach, then the Japanese for strolling.
\end{quote}

\subsection*{You'll want to know: \ja{うろうろ}}
\textbf{\ja{うろうろ}} — \emph{urōro} — \textquotedblleft{}wandering restlessly\textquotedblright{}.

\begin{quote}
\textbf{\ja{えきの} \ja{まえで} \ja{うろうろしていました。}} — \emph{eki no mae de urōro shite imashita} — \textquotedblleft{}I was hanging about in front of the station.\textquotedblright{}
\end{quote}

\begin{cousinweb}
\textbf{\ja{うろうろ}} (\emph{urōro}) is a native mimetic word that pictures restless, back-and-forth wandering: an imitative Japanese form, inherited rather than borrowed, whose older origin is not settled (JMdict).
\end{cousinweb}

\begin{grammarlens}[title={What you've built}]
One more word for getting around.
\end{grammarlens}

\subsection*{Your turn}
\begin{itemize}
\raggedright
  \item \emph{Say it:} \emph{urōro}
  \item \emph{Say it:} \emph{urōro}, once more
  \item \emph{Say it:} \emph{eki no mae de urōro shite imashita} — I was hanging about in front of the station.
  \item \emph{Recall:} say the Japanese for to finally reach, then the Japanese for strolling, then say \emph{urōro} again
\end{itemize}

\begin{tcolorbox}[breakable,colback=teal!4,colframe=teal!35!black,title={Before you move on}]
What does \ja{うろうろ} mean? (\textquotedblleft{}Wandering restlessly\textquotedblright{}.) Say it once more.
\end{tcolorbox}
